\documentclass[10pt,a4paper]{amsart}
\usepackage[margin=19mm]{geometry}
\usepackage{amsmath,amssymb,mathtools}
\usepackage[scr=rsfs,cal=euler]{mathalfa}
\usepackage{xfrac}
\usepackage[unicode,hidelinks,bookmarks=false]{hyperref}
\newcommand{\ie}{i.e.\ }
\newcommand{\cf}{cf.\ }
\newcommand{\resp}{respectively\ }
\newcommand{\Subsection}{\subsection}
\providecommand{\nobreakdash}{\nobreak\hbox{-}}
\setlength{\emergencystretch}{2em}
\multlinegap0pt
\usepackage{bm}
\usepackage{bbm}
\usepackage{dsfont}
\usepackage{xspace}
\usepackage{cancel}
\usepackage{mathtools}
\usepackage{cleveref}
\usepackage{amsthm}
\makeatletter
\@ifundefined{fact}{\newtheorem{fact}{Fact}}{}
\@ifundefined{prop}{\newtheorem{prop}{Proposition}}{}
\makeatother
\renewcommand{\max}{\operatornamewithlimits{max}}
\title[Spectral Methods in Microeconomics]{Spectral Methods in Microeconomics}
\author{Benjamin Golub}
\date{Source: arXiv:2502.12309v2 (CC BY 4.0)}
\begin{document}
\maketitle

\noindent\emph{Source and licence.} Spectral Methods in Microeconomics. Version arXiv:2502.12309v2 (CC BY 4.0), distributed under the Creative Commons Attribution 4.0 International licence, as recorded in the arXiv licence metadata for this exact version and shown on the abstract page https://arxiv.org/abs/2502.12309v2. Full rights notice: licenses/arxiv-2502.12309-CC-BY-4.0.txt.\par\smallskip

\noindent\textit{Editorial scope.} Counted unit: the Prop whose source label is \texttt{thm:poa} in the article (source0.tex lines 489--495, source label \texttt{thm:poa}), delivered together with the article's own complete original proof of it (source0.tex lines 497--514, one \texttt{proof} environment of 18 lines). No mathematics is written, repaired or re-derived: statement and proof are the article's own text line for line. Required context retained uncounted: eq:network\_game\_payoff (lines 400--403); fact:Nash (lines 419--424). Every definition, display and label of those lines is retained unchanged. Omitted from this package and not redistributed: 1--399, 404--418, 425--488, 496--496, 515--1063. Nothing inside a retained excerpt is omitted or altered: every retained body is delivered exactly as the article's lines are written, so no omission receipt of any kind accompanies this package. Citations: the retained excerpts carry no citation command at all, so no bibliography entry of the article is transcribed and no reference is redistributed. Reference adaptation: none was needed and none was made; every reference inside the delivered unit resolves inside it, so no unresolved reference or citation is produced. Printed slips kept literal: no printed word, symbol, hypothesis or number of the counted statement or of its proof was altered. Layout and class: the article's own class is \texttt{notices} with packages including graphicx, subcaption, adjustbox, multicol, adjmulticol, caption, mathbbol, tkz-graph, tikz, float, sectsty, amsrefs, titlesec, amsthm; the wrapper is the lane's pinned \texttt{amsart} editorial wrapper at 10pt on A4, which restates the theorem-like environments the retained text needs and adds the article's own macro definitions that the retained text needs (\texttt{\textbackslash max}), with extra wrapper packages (bm, bbm, dsfont, xspace, cancel, mathtools, cleveref). The delivered numbering is the unit's own. Assets: no retained span contains a figure environment, a graphics-inclusion command, a PDF-page inclusion, a table or a TikZ picture; no image file is copied into this package, the package declares no asset dependency and no image file is redistributed with it. Licence: version arXiv:2502.12309v2 of this article is distributed under the Creative Commons Attribution 4.0 International licence, as recorded in the arXiv licence metadata for this exact version and shown on the abstract page https://arxiv.org/abs/2502.12309v2. A full rights notice is delivered at licenses/arxiv-2502.12309-CC-BY-4.0.txt. Characters: every retained excerpt is pure ASCII, so the delivered engine, XeLaTeX with its default Latin Modern text font, reports no missing character.
\begin{equation}
u_i(x_1, \ldots, x_n) = -\frac{1}{2} \gamma_i x_i^2 +  \left( \beta_i + \sum_{j \neq i} G_{ij} x_j \right)x_i.
\label{eq:network_game_payoff}
\end{equation}

\begin{fact} \label{fact:Nash}
If $\rho(M) < 1$, then there exists a unique Nash equilibrium given by:
\begin{equation} \label{eq:Nash}
x^* = (I - M)^{-1} b.
\end{equation}
\end{fact}

\begin{prop}
\label{thm:poa}
Assume $2\rho(M)<1$, where $\rho(M)$ denotes the spectral radius of $M$. Then,
$
\text{PoA} = \frac{(1-\rho(M))^2}{1-2\rho(M)}.
$
\end{prop}

\begin{proof}[Sketch]

Diagonalize $M$ as $M = W \Lambda W^\top$, where $\Lambda = \text{diag}(\lambda_1, \ldots, \lambda_n)$ and $W$ is an orthogonal matrix whose columns are eigenvectors of $M$. Let $\tilde{x} = W^\top x$ and $\tilde{b} = W^\top b$. The Nash equilibrium and socially efficient actions can be written as
$
\tilde{x}_\ell^* = \frac{\tilde{b}_\ell}{1 -  \lambda_\ell}, \quad \tilde{x}_\ell^{\mathrm{eff}} = \frac{\tilde{b}_\ell}{1 - 2 \lambda_\ell}.
$ The first follows from rewriting \Cref{fact:Nash} in the diagonal basis, and the second follows by solving a very similar system of equations corresponding to the first-order conditions for maximizing $V(x)$.

The total welfare at the Nash equilibrium is
\begin{equation}
V(x^*) = \frac{1}{2} \sum_{\ell=1}^n \frac{\tilde{b}_\ell^2}{(1 -  \lambda_\ell)^2}.  \end{equation} This is obtained by plugging in the condition that all agents are best-responding, $x^*=Mx^* + b$, into agents' utility functions from \cref{eq:network_game_payoff}, noting that in the present case $\beta=b$. This shows that $u_i = \frac{1}{2}(x_i^*)^2$, so that utilitarian welfare is $\frac{1}{2} (x^*)^\top x^*$. We then plug in the above characterization $\tilde{x}_\ell^* = \frac{\tilde{b}_\ell}{1 -  \lambda_\ell}$ in the diagonalized basis. By a similar calculation,
\begin{equation}
V(x^{\mathrm{eff}}) = \frac{1}{2}\, b^\top (I-2M)^{-1} b
 \;=\; \frac{1}{2} \sum_{\ell=1}^n \frac{\tilde{b}_\ell^2}{1 - 2 \lambda_\ell}.
\end{equation}

Thus, when $b$ is proportional to an eigenvector of $M$ with eigenvalue $\lambda$, the welfare ratio is
$\frac{V(x^{\mathrm{eff}})}{V(x^*)}=\frac{(1-\lambda)^2}{1-2\lambda}$, which is increasing in $\lambda$ on $[0,1/2)$. The PoA is therefore maximized when all weight is placed on the eigenvector corresponding to $\lambda_{\max} = \rho(M)$, yielding the result.
\end{proof}

\end{document}
